\chapter{System Design and Production Architecture}
\label{ch:architecture}

The algorithms are settled. This chapter is about everything that makes them survive contact with a
real machine: how the code is split, why there is a cache, what happens when the camera is missing,
how the whole thing is tested without hardware, and how a text overlay is drawn so that it stays
readable.

\section{Architectural separation}
\label{sec:split}

The project is two modules, and the boundary between them is the most important design decision in
the codebase.

\begin{figure}[h]
\centering
\begin{tikzpicture}[
  mod/.style={draw=navy, fill=slatebg, rounded corners=2mm, align=center,
              inner sep=4mm, font=\small\sffamily, text=navy, line width=1pt},
  itm/.style={font=\scriptsize\sffamily, text=slatetext, align=left},
  ar/.style={-{Stealth[length=3mm]}, draw=cyanx, line width=1.4pt},
  arred/.style={-{Stealth[length=3mm]}, draw=redx, line width=1.6pt, dashed}
]
  % hand_tracker
  \node[mod, text width=6.4cm] (t) at (0,0)
    {\textbf{hand\_tracker.py}\\[2pt]
     \rule{6cm}{0.4pt}\\[3pt]
     \begin{tabular}{@{}l@{}}
       camera I/O \textbullet\ \code{cv2.imshow}\\
       BGR$\to$RGB conversion\\
       MediaPipe model ownership\\
       HUD rendering \textbullet\ CLI \textbullet\ logging\\
       fallback chains \textbullet\ teardown
     \end{tabular}};

  % hand_logic
  \node[mod, text width=6.4cm, draw=tealx] (l) at (8.6,0)
    {\textbf{hand\_logic.py}\\[2pt]
     \rule{6cm}{0.4pt}\\[3pt]
     \begin{tabular}{@{}l@{}}
       landmark parsing \textbullet\ cache building\\
       finger \& thumb rules\\
       palm normal \textbullet\ bounding box\\
       EMA \textbullet\ gesture classification\\
       debouncing
     \end{tabular}};

  % the dependency arrow
  \draw[ar] (t.east) -- node[above, font=\scriptsize\sffamily, text=cyanx]{imports} (l.west);

  % what is forbidden
  \draw[arred] (l.south) -- ++(0,-1.5) node[below, font=\scriptsize\sffamily, text=redx, align=center]
       {no \code{import cv2}\\ no \code{import mediapipe}\\ no camera, no disk, no clock};

  % what is gained
  \node[font=\scriptsize\sffamily, text=tealx, align=center] at (8.6,-2.5)
       {$\Rightarrow$ testable in a bare Python environment};

  % hardware
  \node[mod, text width=3.4cm, draw=slatetext, fill=white] (hw) at (-5.6,0)
    {\textbf{hardware}\\[2pt]\rule{3cm}{0.4pt}\\[3pt]
     \begin{tabular}{@{}l@{}}
       webcam device\\
       display server
     \end{tabular}};
  \draw[ar] (hw.east) -- (t.west);
\end{tikzpicture}
\caption{The dependency rule. All hardware, all third-party libraries and all I/O live on the left.
The right-hand module is pure: it can be imported, exercised and exhaustively tested in a bare Python
interpreter with no camera, no display and no OpenCV installed. The dashed red arrow states what must
never appear in it.}
\label{fig:architecture}
\end{figure}

\subsection{Why the split exists}

The original specification asked for a single \code{hand\_tracker.py} containing everything, with
``unit-testable helpers''. Those two requirements contradict each other: a helper that lives in a
file which imports \code{cv2} and \code{mediapipe} cannot be tested without installing both, and
importing the module to test one geometry function drags in a 40-megabyte native library and, on a
headless machine, may fail outright.

Splitting resolves the contradiction cleanly:

\begin{itemize}
  \item \code{hand\_logic.py} imports \code{math} and \code{collections} only. Its functions take
        plain numbers and lists and return plain numbers and lists.
  \item \code{hand\_tracker.py} owns everything that touches the outside world.
\end{itemize}

\begin{keyidea}
\textbf{The rule, stated as a test you can apply to any new line of code:} \emph{would this line work
if there were no camera, no screen and no MediaPipe?} If yes, it belongs in \code{hand\_logic.py}. If
no, it belongs in \code{hand\_tracker.py}.

Everything the project knows about \emph{hands} is on the pure side. Everything it knows about
\emph{computers} is on the other.
\end{keyidea}

\subsection{What the split buys, concretely}

\begin{center}
\rowcolors{2}{slatebg}{white}
\begin{tabularx}{\textwidth}{@{}l X@{}}
\toprule
\rowcolor{white}
\textbf{Benefit} & \textbf{Consequence in this project} \\
\midrule
Fast tests & The logic suite runs in under a second, with no model loading and no hardware. \\
CI without a GPU & The continuous-integration job needs \code{pip install pytest} and nothing else
                  --- no 200 MB wheel, no display server. \\
Deterministic tests & Synthetic coordinates produce exactly predictable answers; a real webcam never
                      would. \\
Portability & The geometry is valid on any platform, at any resolution, for any camera. \\
Debuggability & When the count is wrong you can ask ``is the geometry wrong, or is the input wrong?''
                and answer it by feeding the geometry a known-good point set. \\
\bottomrule
\end{tabularx}
\end{center}

The test suite makes the payoff visible: the majority of the tests construct arrays of coordinates by
hand and assert on booleans --- no hardware anywhere in sight.

\section{Single-pass telemetry caching}
\label{sec:cache}

\subsection{The problem: repeated work in a hot loop}

A naive implementation converts normalized landmarks to pixels wherever it needs them. In one frame,
the code needs pixel coordinates for: the index-tip marker, all five finger tests (six landmarks),
the thumb rule (three landmarks), the palm normal (three landmarks), and the bounding box (all 21).
Recomputing inside each of those means iterating the landmark list five or more times per hand, per
frame, at 30 frames per second --- thousands of redundant multiplications per second.

That is not a catastrophe. It is waste, and waste in a real-time loop is what separates a program
that holds 30 FPS on a modest laptop from one that drops frames under load.

\subsection{The solution: one parse, one dictionary per hand}

\begin{codecard}
\begin{lstlisting}
hands.append(
    {
        "label": label,
        "norm": norm,
        "px": px,
        "bbox": bbox_from_pixels(px),
        "palm_z": z,
        "orientation": orientation_from_palm_z(z),
    }
)
\end{lstlisting}
\end{codecard}

Everything downstream reads from this structure. The cost of the conversion is paid exactly once per
hand per frame.

\begin{center}
\rowcolors{2}{slatebg}{white}
\begin{tabularx}{\textwidth}{@{}l l X@{}}
\toprule
\rowcolor{white}
\textbf{Key} & \textbf{Type} & \textbf{What it stores, and who consumes it} \\
\midrule
\code{label} & \code{str} & \code{"Left"}, \code{"Right"} or \code{"Unknown"}, after any
  \code{swap\_handedness}. Shown in the HUD; used as the state key for smoothing and debouncing. \\
\code{norm} & \code{list[tuple]} & The raw $(x, y, z)$ triples in $[0,1]$. Kept because it is the
  model's own output --- useful for debugging, and free to retain. \\
\code{px} & \code{list[tuple]} & Integer pixel landmarks. Consumed by every geometric rule and by
  the tip marker. \\
\code{bbox} & \code{tuple} & $(x_{\min}, y_{\min}, x_{\max}, y_{\max})$. The hand's screen extent;
  used for labelling and diagnostics. \\
\code{palm\_z} & \code{float} & The signed palm-normal value from Chapter~\ref{ch:math}. \\
\code{orientation} & \code{str} & \code{"Palm"} or \code{"Back"} --- the label derived once from
  \code{palm\_z}, so no consumer has to repeat the sign test. \\
\bottomrule
\end{tabularx}
\end{center}

\begin{edgetrap}{The cache has a lifetime, and it is one frame}
\code{find\_positions} is called with a frame argument, but it never reads that argument --- the
pixels were computed inside the preceding \code{find\_hands} call, using \emph{that} frame's
dimensions. The parameter exists only to preserve the signature the specification asked for.

The consequence is a real constraint: \textbf{the cached pixel coordinates are only valid for a frame
of the size they were computed from}. Reuse the tracker against a differently sized frame without
calling \code{find\_hands} first and the coordinates will be scaled wrongly --- the dot will be drawn
in the wrong place, with no error raised.

The rule to follow: one \code{find\_hands} per frame, then read. Never cache across frames, never mix
sizes. This is documented in the method's docstring precisely because the failure is silent.
\end{edgetrap}

\subsection{Invalid input, and the safe result}

Parsing is defensive because the model's output shape can vary: a hand with fewer than 21 landmarks
cannot be reasoned about, and a missing handedness list is normal in some configurations.

\begin{codecard}
\begin{lstlisting}
hand_landmarks = getattr(results, "multi_hand_landmarks", None)
if not hand_landmarks:
    return []

handedness = getattr(results, "multi_handedness", None) or []
...
landmarks = getattr(hand, "landmark", None)
if not landmarks or len(landmarks) < EXPECTED_LANDMARKS:
    continue
\end{lstlisting}
\end{codecard}

Three defensive decisions, each with a reason:

\begin{itemize}
  \item \code{getattr(..., None)} rather than \code{results.multi\_hand\_landmarks}: a result object
        with no hands may legitimately lack the attribute. \code{getattr} returns the default instead
        of raising \code{AttributeError}.
  \item \code{... or []} handles the ``attribute exists but is \code{None}'' case, which would
        otherwise break \code{len()}.
  \item \code{continue} rather than \code{break} for a short landmark list: one malformed hand does
        not discard the other, valid hand in the same frame.
\end{itemize}

The same defensiveness applies to the finger rules, which return \code{[False] * 5} --- never raise,
never return \code{None} --- when handed nothing usable. That contract is asserted by a test, so
callers may rely on it.

\section{Camera fallback and hardware resiliency}
\label{sec:camera}

Opening a webcam is the least reliable operation in the entire program. The device may be absent, in
use by a video call, or present but only on a different index. A single-shot
\code{cv2.VideoCapture(0)} fails on all three.

\begin{figure}[h]
\centering
\begin{tikzpicture}[
  n/.style={draw=slateline, fill=slatebg, rounded corners=1.2mm, align=center,
            inner sep=2.4mm, font=\scriptsize\sffamily, text=navy, text width=2.7cm},
  a/.style={-{Stealth[length=2mm]}, draw=cyanx, line width=0.9pt},
  f/.style={-{Stealth[length=2mm]}, draw=redx, line width=0.9pt},
  y/.style={-{Stealth[length=2mm]}, draw=tealx, line width=1.2pt}
]
  \node[n] (start) {requested index\\\code{--camera N}};
  \node[n, right=0.9cm of start] (be) {for each backend:\\DSHOW $\to$ MSMF $\to$ ANY};
  \node[n, below=0.7cm of be] (next) {next candidate index\\(1, then 0)};
  \node[n, right=0.9cm of be] (ok) {\textbf{opened}\\set 640$\times$480\\log the index};
  \node[n, below=0.7cm of ok, draw=redx, fill=ambersoft] (fail)
      {\textbf{all candidates failed}\\log an actionable error\\exit code 1};

  \draw[a] (start) -- (be);
  \draw[a] (be) -- node[above, font=\scriptsize\sffamily, text=tealx]{isOpened} (ok);
  \draw[f] (be) -- node[right, font=\scriptsize\sffamily, text=redx, align=left]{not opened} (next);
  \draw[a] (next) -| (be);
  \draw[f] (next) -- (fail);
\end{tikzpicture}
\caption{The camera fallback chain. Both axes are searched --- backends within one index, and then
other indices --- before giving up with a message that names the likely cause.}
\label{fig:camera}
\end{figure}

\begin{codecard}
\begin{lstlisting}
for candidate in [index] + [other for other in (1, 0) if other != index]:
    for backend in backends:
        capture = cv2.VideoCapture(candidate, backend)
        if capture.isOpened():
            capture.set(cv2.CAP_PROP_FRAME_WIDTH, width)
            capture.set(cv2.CAP_PROP_FRAME_HEIGHT, height)
            LOG.info("Camera %d opened at %dx%d.", ...)
            return capture
        capture.release()
    LOG.warning("Camera %d did not open.", candidate)

LOG.error("Could not open a webcam. Check it is connected and not in use by another app.")
return None
\end{lstlisting}
\end{codecard}

\begin{itemize}
  \item \textbf{Backends differ per platform.} On Windows the list is
        \code{[CAP\_DSHOW, CAP\_MSMF, CAP\_ANY]}: DirectShow is the most permissive for consumer
        webcams, Media Foundation is the modern stack, and \code{CAP\_ANY} lets OpenCV choose.
        Elsewhere the list is just \code{[CAP\_ANY]}, because the alternatives do not exist.
  \item \textbf{The candidate list is} \code{[index] + [1, 0] with the requested index removed},
        so \code{--camera 0} tries 0 then 1, while \code{--camera 5} tries 5, then 1, then 0.
  \item \code{capture.release()} on every failure is not optional. Without it, each failed attempt
        leaks a device handle, and on Windows a leaked handle can itself prevent the next attempt
        from succeeding --- the fallback chain would sabotage itself.
  \item \textbf{Width and height are requested, not guaranteed.} \code{cap.set} is a request to the
        driver; the log line reports what was actually granted, which is why it re-reads the
        properties rather than echoing the requested numbers.
\end{itemize}

\begin{edgetrap}{``Could not open a webcam'' --- the four real causes}
The error text is deliberately specific because the causes are mundane and diagnosable:

\begin{enumerate}
  \item The camera is in use by another application (a video call, a browser tab holding the device,
        or a previous run of this program that was killed without cleanup).
  \item The index is wrong --- common on laptops with both an integrated and an external camera, and
        on desktops where a capture card occupies index 0.
  \item Operating-system privacy settings are blocking camera access for desktop applications.
  \item The device is genuinely absent or disabled in Device Manager.
\end{enumerate}

Cause 1 is the most common by a wide margin, and it is the reason teardown matters
(Section~\ref{sec:teardown}).
\end{edgetrap}

\section{Headless mode: verifying without a webcam}
\label{sec:headless}

Every test in the project's logic suite passes without a camera --- but until this feature existed,
\emph{the drawing and rendering path} could only be exercised by a human sitting in front of a
webcam. That is a large untested surface: the mesh drawing, the tip marker, the HUD blending, the
CLI, the image encoding.

\code{--image PATH} closes that gap. It runs the entire pipeline --- load, detect, parse, compute,
draw, write --- on a single still.

\begin{center}
\rowcolors{2}{slatebg}{white}
\begin{tabularx}{\textwidth}{@{}r l X@{}}
\toprule
\rowcolor{white}
\textbf{Step} & \textbf{Call} & \textbf{Purpose} \\
\midrule
1 & \code{cv2.imread(path)} & Returns \code{None} on failure; the code checks and exits non-zero. \\
2 & \code{cv2.resize} & Only if \code{--width} \emph{and} \code{--height} were supplied. A still has
    a native resolution worth preserving. \\
3 & \code{HandTracker(mode=True)} & \code{static\_image\_mode} --- there is no next frame to track
    into. \\
4 & \code{find\_hands} $\to$ \code{annotate\_hands} $\to$ \code{draw\_hud} & \emph{The same
    functions the live loop calls.} No parallel implementation exists to drift out of sync. \\
5 & \code{cv2.imwrite} & Writes \code{<stem>\_out.png} next to the input, then exits 0. \\
\bottomrule
\end{tabularx}
\end{center}

Note the deliberate omission: the still is \textbf{not mirrored}. A photograph is not a selfie, and
flipping it would move every landmark the model found to the wrong side of the image.

\subsection{Continuous integration without a display}

\begin{codecard}
\begin{lstlisting}[style=codetextonly]
name: CI
on: [push, pull_request]

jobs:
  logic-tests:
    runs-on: ubuntu-latest
    strategy:
      matrix:
        python-version: ["3.11", "3.12"]
    steps:
      - uses: actions/checkout@v4
      - uses: actions/setup-python@v5
        with: { python-version: "${{ matrix.python-version }}" }
      - run: pip install pytest
      - run: pytest -q

  pipeline-smoke:
    runs-on: ubuntu-latest
    steps:
      - uses: actions/checkout@v4
      - uses: actions/setup-python@v5
        with: { python-version: "3.12" }
      - run: sudo apt-get update && sudo apt-get install -y libgl1 libglib2.0-0
      - run: pip install -r requirements.txt pytest
      - run: pytest -q
      - run: python hand_tracker.py --image tests/fixtures/hand.png
      - run: test -f tests/fixtures/hand_out.png
\end{lstlisting}
\end{codecard}

\begin{edgetrap}{\code{libGL.so.1}: the OpenCV-on-CI gotcha}
The \code{pipeline-smoke} job installs the real dependencies --- including
\code{opencv-contrib-python}, which is \emph{not} the headless build. Those wheels link against
system graphics libraries that a minimal container does not have, and the failure is a stark

\begin{lstlisting}[style=codetextonly]
ImportError: libGL.so.1: cannot open shared object file: No such file or directory
\end{lstlisting}

It looks like a broken install or a code bug. It is neither: it is a missing shared library in the
runner image. Installing \code{libgl1} (and \code{libglib2.0-0}, which the same wheels frequently
need) fixes it in one line. This is the single most commonly hit OpenCV/CI failure, which is why it
is commented in the workflow itself rather than left as folklore.
\end{edgetrap}

\section{The alpha-blended HUD}
\label{sec:hud}

Drawing white text directly onto a webcam image fails the moment anything bright is behind it --- a
window, a lamp, a white wall. The fix is to darken a rectangle behind the text before writing it.

\begin{mathdeepdive}{Alpha blending}
For each pixel in the card region, the output is a weighted average of what was there and the
overlay:

$$I_{\text{out}} = (1-\alpha)\, I_{\text{roi}} + \alpha\, I_{\text{overlay}}$$

The project uses $\alpha = 0.6$ with a black overlay:

$$I_{\text{out}} = 0.4\, I_{\text{roi}} + 0.6 \times 0 = 0.4\, I_{\text{roi}}$$

So a bright background pixel of 255 becomes $102$ --- comfortably dark --- while a dark background
of 40 becomes $16$. The result is that the card is visible as a subtle darkening on any background,
and white text on top is readable at any luminance. Note that the formula is applied
\emph{per channel}, so hues are preserved and only brightness is scaled.
\end{mathdeepdive}

\begin{codecard}
\begin{lstlisting}
card = frame[y0:y1, x0:x1]
frame[y0:y1, x0:x1] = cv2.addWeighted(card, 1.0 - HUD_ALPHA, np.zeros_like(card), HUD_ALPHA, 0)
\end{lstlisting}
\end{codecard}

Two details worth noticing:

\begin{itemize}
  \item \code{frame[y0:y1, x0:x1]} is a \textbf{view}, not a copy. Writing back into it modifies the
        original frame in place, which is why the assignment on the second line is necessary rather
        than optional --- \code{addWeighted} returns a new array.
  \item The ROI is clamped to the image, and the function returns early if the clamp produces an
        empty region:
        \code{if x1 <= x0 or y1 <= y0: return frame}. Without that guard, a frame smaller than the
        card raises on the slice assignment. A test covers this case with a $20\times20$ frame.
\end{itemize}

\begin{edgetrap}{Fixed-size cards overflow with two hands}
The card is sized dynamically --- both dimensions --- from the text it must hold:

\begin{lstlisting}
sizes = [cv2.getTextSize(line, ...)[0] for line in lines]
text_width = max(width for width, _ in sizes)
line_height = text_height + 12
x1 = min(width, x0 + text_width + 2 * HUD_PAD)
y1 = min(height, y0 + 2 * HUD_PAD + line_height * len(lines))
\end{lstlisting}

A hard-coded box would work until the moment two hands are detected, at which point the HUD gains two
extra lines and the text spills out of the card --- over the live video, unreadable. Measuring the
text first makes the card correct for one hand, two hands, or a long ``Rock On'' gesture label.
\end{edgetrap}

\section{Teardown and resource safety}
\label{sec:teardown}

\begin{codecard}
\begin{lstlisting}
try:
    while True:
        ok, frame = capture.read()
        ...
finally:
    capture.release()
    tracker.close()
    cv2.destroyAllWindows()
\end{lstlisting}
\end{codecard}

The \code{finally} block runs on every path out of the loop: the user pressing \code{q}, an
\code{Esc}, a failed frame read that breaks the loop, or an exception. Each call releases a distinct
resource:

\begin{itemize}
  \item \code{capture.release()} returns the camera device to the operating system. Skipping it is
        what produces the ``camera already in use'' error on the next run.
  \item \code{tracker.close()} releases the MediaPipe graph and its native resources.
  \item \code{cv2.destroyAllWindows()} removes the GUI window. Without it, a crashed run can leave an
        orphaned window that keeps a stale frame on screen and is not obviously connected to any
        process.
\end{itemize}

\begin{keyidea}
\textbf{The general principle:} acquire resources in \code{\_\_init\_\_} or at the start of a
function, and release them in \code{finally} --- never at the end of the happy path only. The happy
path is the one case that is already working.
\end{keyidea}

\section{The full data flow}

\begin{figure}[h]
\centering
\begin{tikzpicture}[
  b/.style={draw=slateline, fill=slatebg, rounded corners=1mm, align=center,
            inner sep=1.6mm, font=\scriptsize\sffamily, text=navy,
            text width=2.4cm, minimum height=1.05cm},
  io/.style={b, draw=amberx, fill=ambersoft},
  pure/.style={b, draw=tealx, fill=tealsoft},
  ar/.style={-{Stealth[length=1.8mm]}, draw=cyanx, line width=0.8pt},
  back/.style={-{Stealth[length=1.8mm]}, draw=slatetext, line width=0.7pt, dashed}
]
  % ---- top row: left to right ----
  \node[io]   (cam)  at ( 1.30, 1.40) {camera frame\\BGR $480\times640\times3$};
  \node[b]    (flip) at ( 4.25, 1.40) {\code{cv2.flip}\\mirror};
  \node[b]    (rgb)  at ( 7.20, 1.40) {\code{cvtColor}\\BGR$\to$RGB};
  \node[b]    (palm) at (10.15, 1.40) {BlazePalm\\palm box};
  \node[b]    (lm)   at (13.10, 1.40) {landmark\\regression};

  % ---- bottom row: right to left ----
  \node[pure] (parse) at (13.10,-0.60) {\code{parse\_hand\_data}\\one-pass cache};
  \node[pure] (geom)  at (10.15,-0.60) {\code{fingers\_state}\\+ palm normal};
  \node[pure] (db)    at ( 7.20,-0.60) {\code{GestureDebouncer}};
  \node[b]    (hud)   at ( 4.25,-0.60) {alpha-blended\\HUD};
  \node[io]   (scr)   at ( 1.30,-0.60) {window\\(\code{imshow})};

  % ---- flow ----
  \draw[ar] (cam)  -- (flip);
  \draw[ar] (flip) -- (rgb);
  \draw[ar] (rgb)  -- (palm);
  \draw[ar] (palm) -- (lm);
  \draw[ar] (lm)   -- (parse);
  \draw[ar] (parse)-- (geom);
  \draw[ar] (geom) -- (db);
  \draw[ar] (db)   -- (hud);
  \draw[ar] (hud)  -- (scr);

  % ---- feedback loop, left-hand side ----
  \draw[back] (0.10,-0.60) -- (-0.80,-0.60) -- (-0.80,1.40) -- (0.10,1.40);
  \node[font=\scriptsize\sffamily, text=slatetext, rotate=90, anchor=south]
       at (-0.98, 0.40) {next frame, 30$\times$/s};

  % ---- legend ----
  \node[draw=amberx, fill=ambersoft, minimum size=3.4mm, inner sep=0pt] at (4.30,-1.90) {};
  \node[font=\scriptsize\sffamily, text=slatetext, anchor=west] at (4.62,-1.90)
       {I/O and the model (\code{hand\_tracker.py})};
  \node[draw=tealx, fill=tealsoft, minimum size=3.4mm, inner sep=0pt] at (4.30,-2.30) {};
  \node[font=\scriptsize\sffamily, text=slatetext, anchor=west] at (4.62,-2.30)
       {pure logic (\code{hand\_logic.py})};
\end{tikzpicture}
\caption{One frame, end to end. The trace runs left to right along the top row --- mirror, colour
conversion, detect, regress --- then drops into the pure-logic row and travels back right to left to
the window. The dashed line on the left is the only thing making this a video rather than a snapshot.}
\label{fig:pipeline}
\end{figure}

\begin{keyidea}
\textbf{Chapter summary.} The codebase is split along exactly one seam --- pure logic versus
everything that touches hardware --- and that seam is what makes the project testable at all. Each
frame converts its landmarks once into a per-hand dictionary, which every downstream rule reads.
Opening a camera searches both backends and indices before failing with a message that names the
likely cause, and all resources are released in a \code{finally} block. The \code{--image} mode runs
the identical pipeline on a still, which is what allows the whole program, drawing included, to be
verified in CI without a display server.
\end{keyidea}
